\mediumtitle{Teaching activity}

\bigskip

\smalltitle{To BSc and MSc students}

\begin{doubletablelist}
	
	\stlist{a.y. 2025-26 - present}{\textbf{Co-Lecturer} for the \emph{``Multimorbidity''} elective teaching activities for the \emph{Faculty of Medicine and Surgery - MedInTO}, Department of Clinical and Biological Science,  University of Turin. The session led by Alessia Visconti (2h) explains how computational medicine is used to predict the clinical trajectories and to stratify patients with multimorbidity. Language: English. \\
	& \textbf{Lecturer} for the \emph{``Introduction to Workflow Orchestration Managers''} workshop for the \emph{Advanced Bioinformatics} module, MSc in Genomic Medicine, King’s College London. The workshop (3h) follows the ``live coding'' principle Alessia Visconti champions and offers a primer on Nextflow, a very popular workflow orchestration manager. Language: English.
	}
	
	\stlist{a.y. 2024-25 - present}{\textbf{Lecturer} for the \emph{``Medical statistics''} module within the \emph{``Evidence-based nursing''} course, Nursing School, Department of Clinical and Biological Science, University of Turin. The module (15h) covers the fundamentals of biostatistics and is run in four parallel sessions in two different university campuses, Cuneo (from a.y. 2024-25) and Orbassano (from a.y. 2026-27).  Language: Italian.\\
		& \textbf{Lecturer} for the \emph{``Medical statistics''} module within the \emph{``Preparatory to research''} course, Degree in Psychiatric Rehabilitation Techniques, Department of Clinical and Biological Science, University of Turin. The module (24h) covers the fundamentals of biostatistics through lectures and practical activities. Language: Italian.	\\	
		& \textbf{Co-lecturer} for the \emph{``Leveraging Large Language Models for hands-on learning in medical education and research''} elective teaching activities for the \emph{Faculty of Medicine and Surgery - MedInTO}, Department of Clinical and Biological Science,  University of Turin. The activity (10h) is organised as a set of lectures, hands-on activities, and discussions. It covers how large language models work, their hidden costs, and whether to use them in reading, writing, studying, and in doing clinical research and work. Language: English. 
		}
	
	\stlist{a.y. 2024-25}{\textbf{Co-lecturer} for the \emph{``Methodology of social-health research''} module within the \emph{``Sociology and Methodology of Health and Social Care Research''} course, Degree in Professional Education, University of Turin. The module (20h) covered the fundamentals of biostatistics. Language: Italian.
	}
	
	\stlist{a.y. 2013-14}{\textbf{Teaching assistant} for the \emph{``Human Molecular Genetics''} MSc Department of Genomics of Common Diseases, Imperial College London. Alessia Visconti supported students during the practical sessions within the following courses:\\
		%& \hskip0.5cm Workshops on: \\
		& \hskip1cm - The Unix Shell\\
		& \hskip1cm - R Programming\\
		& \hskip1cm - Exploratory Data Analysis and Probability\\
		& \hskip1cm - Quantitative genetics\\
		& \hskip1cm - Next Generation Sequencing Data Analysis
	}
				
	\stlist{a.y. 2013-14}{\textbf{Teaching assistant} for the \emph{``Data analysis''} course, Department of Biological Science, University of Turin. Alessia Visconti supported students during the practical sessions on R programming (10h) as well as during one-to-one meetings. Language: Italian. 
	}
	
	\stlist{a.y. 2012-13 \& 2013-14}{\textbf{Teaching assistant} for the \emph{``Operating System''} course, Department of Computer Science, University of Turin. Alessia Visconti supported students during the practical sessions (25h) as well as during one-to-one meetings. Language: Italian.
	}
	
	\stlist{a.y. 2012-13}{\textbf{Lecturer} for the \emph{``Operating System and Networking''} course, Interfaculty School of Strategic Studies, University of Turin. Alessia Visconti was responsible for the practical sessions covering the basis of GNU/Linux, the Unix shell, and process management (28h). She designed the final project, \emph{i.e.}, the development of a basic client/server application in C. Language: Italian.
	}
	
	\stlist{a.y. 2010-11 \& 2011-12}{\textbf{Teaching assistant} for the \emph{``Database''} course, Department of Computer Science, University of Turin. Alessia Visconti supported students during the practical sessions (25h) as well as during one-to-one meetings. Language: Italian.
	}

			
	\stlist{a.y.  2010-11}{\textbf{Teaching assistant} for the \emph{``Formal Language''} course, Department of Computer Science, University of Turin. Alessia Visconti supported students during the practical sessions (25h) as well as during one-to-one meetings. Language: Italian.\\
		&\textbf{Teaching assistant} for the \emph{``Statistics and data mining with SAS''} course, Department of Mathematics, University of Turin. Alessia Visconti prepared all the course online material, including a self-evaluation questionnaire, and recorded 3 hours of video lessons on SAS Enterprise Miner. Language: Italian.
	}
		
	\stlist{a.y.  2009-10}{\textbf{Lecturer} for the \emph{``Computer Science''} course, Department of Letters and Philosophy, University of Turin. The hands-on course (30h in two parallel sessions) covered the MS Office suite. Language: Italian.
	}
	
	\stlist{a.y. 2005-06 \& 2006-07}{\textbf{Teaching assistant} for the \emph{``Program Languages - JAVA''} course, Department of Computer Science, University of Turin. Alessia Visconti supported students during the practical sessions. Language: Italian.
	}
	
	\stlist{a.y. 2004-05}{\textbf{Teaching assistant}  for the \emph{``Program Languages - C''} course, Department of Computer Science, University of Turin. Alessia Visconti supported students during the practical sessions. Language: Italian.
	}

\end{doubletablelist}


\smalltitle{To medical residents}

\begin{doubletablelist}
	
	\stlist{a.y. 2025-26 -- present}{ \textbf{Co-lecturer} for the \emph{``Introduction to biostatistics''} course for medical residents at University of Turin. The course (7h) covers the fundamentals of biostatistics. Language: Italian.
		}

	\stlist{a.y. 2024-25 -- present}{ \textbf{Co-lecturer} for the \emph{``Integrating Large Language Models into clinical work and research: a critical approach''} course for medical residents at University of Turin. The course (7h) is organised as a set of lectures, hands-on activities, and discussions. It covers how large language models work, their hidden costs, and whether to use them in the clinic and in clinical research. Language: English.
		}

\end{doubletablelist}



\smalltitle{To postgraduate and PhD students, and postdoctoral researchers}

\begin{doubletablelist}
	
	\stlist{a.y. 2025-26 -- present}{ \textbf{Lecturer} for the \emph{``Introduction to biostatistics''} workshop for the PhD program in Complex Systems for Quantitative Biomedicine, at the University of Turin.  The workshop (2h) covers, through practical exercises, the fundamentals of biostatistics. Language: Italian. \\
	& \textbf{Co-lecturer} for the \emph{``Machine Learning in Clinical Research: Methods, Interpretability, and Limitations''} workshop for the Advanced Master programme in Epidemiology, Department of Medical Sciences, University of Turin. The workshop (8h) offers a primer on supervised and unsupervised machine learning, and is organised as a set of lectures with a final hands-on activity. Language: Italian.
		}
	
	\stlist{a.y. 2024-25 -- present}{\textbf{Lecturer} for the \emph{``Introduction to Git''} workshop for the PhD program in Complex Systems for Quantitative Biomedicine, at the University of Turin. The workshop (8h) is centred around the idea of ``live coding'', where attendees code along with the instructors thus getting useful hands-on experience and improving their ability to explore the topics on their own. Language: English.\\
		& \textbf{Lecturer} for the \emph{``A Crash Course on Machine Learning for Beginners: Practice''} workshop for the PhD program in Complex Systems for Quantitative Biomedicine, at the University of Turin. The workshop (8h) follows the same ``live coding'' principle, and offers an engaging hands-on session on supervised learning. Language: English.
	}
	
	\stlist{a.y. 2019-20 to 22-23}{\textbf{Instructor} for several programming and data analysis workshops held at King's College London covering the following topics:\\
		& \hskip1cm - The Unix Shell\\
		& \hskip1cm - Version Control and collaboration with Git/GitHub\\
		& \hskip1cm - Python Programming\\
		& \hskip1cm - R Programming\\
		& \hskip1cm - Introduction to Working with Data\\
		& \hskip1cm - OpenRefine\\
		& These workshops were designed for PhD students and early career researchers but were open to researchers at every level, including PIs, and offered regularly (roughly twice a year). They have always been evaluated as ``excellent'' or ``exceptional'' by the attendees. The workshops are centred around the idea of ``live coding''. Language: English.
	}				
	
	% \stlist{2016 - 2023}{\textbf{Co-organiser} of the Regulatory Genomics journal club at the Department of Twins Research \& Genetic Epidemiology, King's College London. The journal club discusses papers on the latest achievements and methods in the fields of regulation of gene expression, epigenetics, splicing, evolution, and related topics. \\
	% 	&
	% }
		
	\stlist{a.y. 2021-22}{\textbf{Instructor} for \emph{``The Unix Shell''} and the \emph{``Version Control and collaboration with Git/GitHub''} workshops (18h) for students and researchers of the MRC Gambia and the AKU Nairobi Research Centres. Both workshops were rated as ``exceptional'' by the attendees.  Language: English. \\
	 	& \textbf{Instructor} for the \emph{``Metagenomics Data Analysis: Investigating the invisible world of microbes''}, Department of Twins Research \& Genetic Epidemiology, King's College London. The entry-level workshop (14h), which included frontal lessons and hands-on sessions, was designed for researchers at every level (including but not limited to PhD students, post-doctoral researchers, and PIs). Language: English.
	}


\end{doubletablelist}

\smalltitle{Outside academia}


\begin{doubletablelist}


\stlist{a.y. 2021-22}{ \textbf{Instructor} for a \emph{``Version Control and collaboration with Git/GitHub''} workshop at the UK Health Security Agency (UKHSA, 9h).}

\end{doubletablelist}

